\section{Universal section bases and the two global section laws}
\label{sec:global-section-laws-v052}

\status{GSL1--GSL26 reconstruct the finite-complex and geometric
arguments supplying the original section laws.  The conclusions concern
arbitrary coefficient rings and the actual line-bundle sections, including
every residue-section lift in \texttt{SecTwistSt}.  The precise upstream
geometric premises are recorded below; they are not silently promoted to
proved global laws.  All source references use the accepted formal proof
at commit \texttt{7cbf7a0c32fca1badf05931879a37c3f8a2b03b4}.}

The purpose of this section is to remove the linear-algebra and
section-transport work from the premises retained in NFL22 and TDR13.
Its central construction is an identity of matrices over the original
ring.  Consequently, neither the coefficient ring nor a later base
change is required to be Noetherian, reduced or flat.  The final
statements preserve the original polarized abelian scheme, its line
bundle, the chosen point and every subsequent affine test map.

Here and below, a line bundle $M$ is the source's actual cocycle
\texttt{LB}: an open cover $U_i$, units $g_{ij}$ on intersections,
and $g_{ij}g_{jk}=g_{ik}$.  A section is an actual family $s_i$ with
$s_i=g_{ij}s_j$.  We write $\Gamma(X,M)$ for this module; this is
\texttt{LB.Sec} with its module structure, or \texttt{SecM} after
the specified \texttt{LPow} chart conversion.  We identify
$\Gamma(\operatorname{Spec}R,\mathcal O)$ with $R$ by the canonical
ring isomorphism, and transport back through that isomorphism whenever
the original statement uses its literal coefficient ring.

For $f:X\to\operatorname{Spec}R$, let $c_f:R\to\Gamma(X,\mathcal O)$
be the structure map.  The source's $\operatorname{pw}(M,t)$ has
transition functions $g_{ij}^{t}$, for any integer $t$; thus negative
powers cause no difficulty.  The source's
$\operatorname{tw}(M,H,t)$ is $M\otimes H^t$, on the intersection
cover.  All isomorphisms used below are represented by actual units on
the appropriate common charts.  An equality of classes in a Picard group
alone is not used as a section comparison.

We abbreviate \texttt{FirstCohomologyVanishes} by $\mathrm{V}(M)$.
Its meaning is the following explicit assertion, for \emph{every}
open cover $V_a\subseteq U_{i(a)}$: if
\[
 \phi_{ab}+g_{i(a)i(b)}\phi_{bd}=\phi_{ad}
 \quad\hbox{on }V_a\cap V_b\cap V_d,
\]
then $\phi_{ab}=\psi_a-g_{i(a)i(b)}\psi_b$ for suitable actual
functions $\psi_a\in\Gamma(V_a,\mathcal O)$.  This notation does
not introduce an extra cohomology comparison theorem.

\subsection{One kernel basis before every base change}

% BEGIN GSL1
\paragraph{GSL1: section coordinates, pullback and uniqueness.}
\label{claim:GSL1}
Compatible functions on any open refinement of a line-bundle cover
determine a unique section.  The construction is additive, commutes with
multiplication by a function from the base, and commutes with pullback.
For a finite linear combination of pulled line bundles, these assertions
hold for the literal \texttt{LinComb} charts and the source's
\texttt{IsPullSecLC}; for a single power they give \texttt{IsPullSec}.

\begin{proof}
Suppose $W_a\subseteq U_{i(a)}$ cover $X$ and the prescribed coordinates
are $v_a$.  On $W_a\cap U_j$, the coordinate in chart $j$ must be
$g_{j,i(a)}v_a$.  The cocycle identity shows that these functions agree
on overlaps.  The sheaf axiom glues them on $U_j$; the same identity gives
the section compatibility between $U_j$ and $U_k$.  This also proves
uniqueness, since equality on the $W_a$ forces equality after changing
to every other chart.  Addition and scalar multiplication commute with
restrictions and multiplication by transition units, so uniqueness gives
the asserted linearity.

If $\pi:X'\to X$, pull back every coordinate and every transition unit.
The functions still satisfy the same compatibility equation, now on
$\pi^{-1}U_i$.  Repeating a pullback gives the pullback by the composite,
because the maps of structure sheaves compose.  For
$\bigotimes_{k\in K}\phi_k^*L^{\epsilon_k}$ the chart indexed by
$j:K\to\operatorname{Charts}(L)$ is
$\bigcap_k\phi_k^{-1}U_{j(k)}$, and its transition is the product
$\prod_k\phi_k^*g_{j(k)j'(k)}^{\epsilon_k}$.  Pullback of this
product is the transition on the corresponding literal common chart.
The coordinate equality just established is exactly
\texttt{IsPullSecLC}.  Taking $K=\{1\}$ gives the power version.
The redundant one-element index in \texttt{LB.pull} is removed by the
mutually inverse chart maps $j\mapsto(\ast\mapsto j)$ and
$J\mapsto J(\ast)$, so no choice of a different section is involved.
\end{proof}
% END GSL1

% BEGIN GSL2
\paragraph{GSL2: a denominator and a left inverse over the original ring.}
\label{claim:GSL2}
Let $R$ be a commutative ring, $K$ a field, $\gamma:R\to K$ a ring
map, and $A:R^l\to R^m$ a matrix.  If $A_K$ is injective, there are
$r\in R$ and a matrix $T:R^m\to R^l$ such that
\[
 \gamma(r)\ne0,\qquad TA=rI_l. \tag{GSL.1}
\]

\begin{proof}
Choose $l$ rows whose restriction $N$ of $A_K$ has nonzero determinant;
they exist because $A_K$ has rank $l$.  Let $J:R^m\to R^l$ select
these same rows over $R$, set $N=JA$, $r=\det N$, and
$T=\operatorname{adj}(N)J$.  The adjugate identity over an arbitrary
commutative ring gives $TA=rI_l$.  Determinants commute with ring maps,
so $\gamma(r)=\det N_K\ne0$.  For $l=0$ use the empty determinant
$r=1$ and the unique empty matrix.  No cancellation in $R$ has been used.
\end{proof}
% END GSL2

% BEGIN GSL3
\paragraph{GSL3: the universal kernel certificate.}
\label{claim:GSL3}
Suppose
\[
 R^a\xrightarrow{D_0}R^b\xrightarrow{D_1}R^c,
 \qquad D_1D_0=0,
\]
is exact at $K^b$ after applying a ring map $\gamma:R\to K$ to a
field.  There are $n\ge0$, $r\in R$, and matrices
$B:R^n\to R^a$, $C:R^a\to R^n$, $H:R^b\to R^a$ with
\[
 \gamma(r)\ne0,\qquad
 D_0B=0,\qquad CB=rI_n,\qquad BC+HD_0=rI_a. \tag{GSL.2}
\]

\begin{proof}
Select columns of $(D_0)_K$ that form a basis of its image, and let
$S:R^k\to R^a$ be their coordinate inclusion.  GSL2 applied to $D_0S$
gives $r_0,T_0$ with $\gamma(r_0)\ne0$ and
$T_0D_0S=r_0I_k$.  The stacked matrix
\[
 A=\begin{pmatrix}T_0\\D_1\end{pmatrix}:R^b\longrightarrow R^{k+c}
\]
is injective over $K$.  Indeed, if $A_Kw=0$, exactness puts $w$ in
$\operatorname{im}(D_0)_K=\operatorname{im}(D_0S)_K$; write
$w=(D_0S)_Kv$.  Then $0=(T_0)_Kw=\gamma(r_0)v$, hence $v=w=0$.
Apply GSL2 again to obtain $T_1A=r_1I_b$, with $\gamma(r_1)\ne0$.

Set $Q=r_0I_a-ST_0D_0$.  Then
\[
 QS=0,\qquad T_0D_0Q=0,\qquad D_1D_0Q=0.
\]
Multiplying the last two equalities by $T_1$ gives $r_1D_0Q=0$;
this is why the second denominator is retained even when $R$ has
zero divisors.  Let $E:R^n\to R^a$ include the complementary
coordinates to those selected by $S$, where $n=a-k$.  Thus
$E^tE=I_n$, $E^tS=0$ and $SS^t+EE^t=I_a$.  Consequently
$QEE^t=Q$ and $E^tQE=r_0I_n$.  Define
\[
 r=r_1r_0,\qquad B=r_1QE,\qquad C=E^t,\qquad H=r_1ST_0.
\]
The first three identities in (GSL.2) now follow directly, and
\[
 BC+HD_0=r_1Q+r_1ST_0D_0=r_1r_0I_a.
\]
Both field images of $r_0,r_1$ are nonzero.  The construction also
covers $k=0$ and $n=0$, using empty matrices throughout.
\end{proof}
% END GSL3

% BEGIN GSL4
\paragraph{GSL4: a kernel basis after every permitted ring map.}
\label{claim:GSL4}
Choose $r,n,B,C,H$ as in GSL3, once and for all.  For every commutative
ring $R'$ and every $\varphi:R\to R'$ for which $\varphi(r)$ is a
unit, the map
\[
 R'^n\longrightarrow\ker(D_0)_{R'},\qquad z\longmapsto B_{R'}z
 \tag{GSL.3}
\]
is an isomorphism, where $\ker(D_0)_{R'}$ means the kernel of the
\emph{base-changed matrix}.  Its inverse on a vector $v$ in that kernel
is $\varphi(r)^{-1}C_{R'}v$.

If $R$ is local, $\kappa$ its residue field, $K/\kappa$ a field
extension and $\gamma:R\to\kappa\hookrightarrow K$, then the
same $r$ is already a unit of $R$.  The basis therefore works over $R$
and $\kappa$, and every vector in $\ker(D_0)_\kappa$ lifts to a
vector in $\ker D_0$.

\begin{proof}
All identities (GSL.2) remain true after applying $\varphi$.  Thus
$D_0B=0$ makes (GSL.3) well defined.  If $D_0v=0$, then
$BCv=\varphi(r)v$, which gives existence of the displayed coordinates.
The identity $CB=\varphi(r)I_n$ gives uniqueness.  These identities
also prove the assertion if $R'$ is the zero ring.  Nothing asserts
that tensoring an arbitrary kernel is exact; the identities prove the
needed conclusion for these particular matrices.

In the local case, $\gamma(r)\ne0$ implies that the residue of $r$
is nonzero, hence $r$ lies outside the maximal ideal and is a unit.
Express a residue vector uniquely as $\overline B\,\overline z$.
Lift its finitely many coordinates along the surjection
$R\to\kappa$, obtaining $z\in R^n$.  Then $Bz$ is a cycle over $R$
and reduces to the prescribed vector.  The basis is chosen before the
residue vector, not anew for each one.
\end{proof}
% END GSL4

\subsection{Finite complexes with the required universal coordinates}

% BEGIN GSL5
\paragraph{GSL5: replacing three finite projectives by three finite matrices.}
\label{claim:GSL5}
Let $P^0\xrightarrow{d_0}P^1\xrightarrow{d_1}P^2$ be a complex
of finite projective $R$-modules.  There are finite free modules
$F^0,F^1,F^2$, matrices $D_0,D_1$ with $D_1D_0=0$, and a fixed
map $p:F^0\to P^0$ such that, after \emph{every} ring map out of $R$,
$p$ induces an isomorphism on the degree-zero kernels.  Exactness at
the middle term of the projective complex implies exactness at the
middle term of the free complex.  The same assertions hold after
tensoring with an arbitrary $R$-module.

\begin{proof}
Choose split presentations
$p_i:E_i=R^{n_i}\to P^i$, $s_i:P^i\to E_i$ with $p_is_i=1$.
Put $F^0=E_0$, $F^1=E_1\oplus E_0$,
$F^2=E_2\oplus E_1\oplus E_0$, and set
\begin{align*}
 D_0(v)&=(s_1d_0p_0v,\ v-s_0p_0v),\\
 D_1(w_1,w_0)&=(s_2d_1p_1w_1,\ w_1-s_1p_1w_1,\ s_0p_0w_0).
\end{align*}
Each component of $D_1D_0$ vanishes, by $d_1d_0=0$ and the
splitting identities.  A vector in $\ker D_0$ satisfies
$v=s_0p_0v$ and $d_0p_0v=0$, so $p_0$ identifies this kernel with
$\ker d_0$, with inverse $s_0$.

Suppose $D_1(w_1,w_0)=0$.  Then $w_1=s_1p_1w_1$,
$d_1p_1w_1=0$, and $p_0w_0=0$ (apply $p_0$ to the third
component).  If the projective complex is exact, choose $x\in P^0$
with $d_0x=p_1w_1$.  The vector $v=s_0x+w_0$ satisfies
$D_0v=(w_1,w_0)$.  Conversely $D_1D_0=0$ gives the other inclusion.
All displayed maps and identities tensor to the same formulas; the
splittings survive arbitrary tensor products.  Finally, if $e_k$ are
the standard basis of $E_0$, the fixed elements $u_k=p_0e_k$
give the kernel comparison by
$v\mapsto\sum_kv_k\otimes u_k$ after base change.  These fixed
coordinates will become the \texttt{Perf.u} functions.
\end{proof}
% END GSL5

% BEGIN GSL6
\paragraph{GSL6: a finite complex over a Noetherian ring.}
\label{claim:GSL6}
Let $R$ be Noetherian and $C^\bullet$ a nonnegative cochain complex,
zero in degrees above $d$, whose cohomology modules are finite.
There exist a complex $K^\bullet$, also zero above $d$, and a
quasi-isomorphism $\phi:K^\bullet\to C^\bullet$, with every $K^p$
finite and every $K^p$ for $p\ge1$ finite free.  No finiteness of the
individual $C^p$ is assumed.

\begin{proof}
We give the finite lifting step underlying a descending construction.
Suppose finite $F^1$, maps $\phi_1:F^1\to C^1$ and
$\delta_1:F^1\to F^2$ have already been chosen.  In
$C^0\oplus F^1$ consider
\[
 N=\{(c,k):d_Cc=\phi_1k,\ \delta_1k=0\}.
\]
Its projection $J$ to $F^1$ is finite because $R$ is Noetherian.
Choose finitely many generators of $J$ and lift them to $N$.  Also
choose cycles in $C^0$ lifting generators of
$\ker d_C/(B\cap\ker d_C)$, where $B$ is the boundary submodule
at the current stage.  This quotient is finite by the assumed
cohomology finiteness.  Take a finite free module on both sets of
chosen lifts, with maps $\phi_0$ and $\delta_0$ given by their first
and second coordinates (the second coordinate of the cycle lifts is
zero).  Then
\[
 d_C\phi_0=\phi_1\delta_0,\quad
 \delta_1\delta_0=0,
\]
and for every $(c,k)\in N$ there is $x$ with
$\delta_0x=k$ and $c-\phi_0x\in B$.  Indeed first match $k$ using
the selected generators of $J$; the remaining first coordinate is a
cycle, and the selected cohomology generators match it modulo $B$.

Start with zero modules above degree $d$ and apply this step downward,
using $B=\operatorname{im}d_C^{p-1}$ at degree $p>0$, and $B=0$
at degree zero.  It produces finite free modules $F^p$, a chain map
$\psi:F^\bullet\to C^\bullet$, and the following lifting
properties:
\begin{enumerate}
 \item If $d_Cc=\psi_1k$ and $\delta_1k=0$ in degree zero,
 there is $x\in F^0$ with $\psi_0x=c$, $\delta_0x=k$.
 \item At degree $p+1$, if $d_Cc=\psi_{p+2}k$ and
 $\delta_{p+2}k=0$, there are $c'\in C^p$, $k'\in F^{p+1}$
 such that $d_Cc'+\psi_{p+1}k'=c$ and $\delta_{p+1}k'=k$.
\end{enumerate}
These are just the preceding lifting step applied successively to the
shifted complex; the boundary error at a positive degree supplies $c'$.

Replace $F^0$ by
$K^0=\operatorname{im}(\psi_0,\delta_0)\subset C^0\oplus F^1$,
and set $K^p=F^p$ for $p\ge1$.  Define $\phi_0$ and $d_K^0$
by the two projections.  The module $K^0$ is finite since it is an
image of $F^0$.  The map $(\phi_0,d_K^0)$ is injective by its
definition; the first lifting property says that its image is the
kernel of the next cone differential.  The second property gives
exactness at all subsequent cone terms.  Explicitly, the resulting
exact sequence is
\[
 0\longrightarrow K^0\xrightarrow{(\phi_0,d_K)}
 C^0\oplus K^1\xrightarrow{(c,k)\mapsto(d_Cc-\phi k,-d_Kk)}
 C^1\oplus K^2\longrightarrow\cdots . \tag{GSL.4}
\]
The same lifting properties show surjectivity and injectivity on each
cohomology module: a cycle lifts modulo a boundary by the second
property with $k=0$, and a cycle mapping to a boundary is a boundary
by that property with the given $k$.  In degree zero the first
property and the initial injection apply.  Thus $\phi$ is a
quasi-isomorphism, as claimed.
\end{proof}
% END GSL6

% BEGIN GSL7
\paragraph{GSL7: flat cones give a universally valid projective model.}
\label{claim:GSL7}
Under GSL6, assume in addition that every $C^p$ is flat over $R$.
Then $K^0$ is finite projective, all $K^p$ are finite projective,
and $\phi$ remains a quasi-isomorphism after tensoring with
\emph{any} $R$-module.  In particular it remains one after every
ring map $R\to R'$.

\begin{proof}
For a short exact sequence $0\to A\to B\to C\to0$ with $B,C$
flat, $A$ is flat and this short exact sequence remains exact after
every tensor product.  To see the latter statement, the only possible
failure is injectivity on the left, measured by
$\operatorname{Tor}_1(C,-)=0$.  The same long exact Tor sequence
gives $\operatorname{Tor}_1(A,-)=0$ since both adjacent higher Tor
groups of $B,C$ vanish.  These are precisely the elementary flatness
criteria used here.

All terms $C^p\oplus K^{p+1}$ in (GSL.4) are flat and the sequence
is bounded on the right.  Starting with its terminal zero image and
working leftward, apply the preceding fact to each short exact sequence
formed by an image, a cone term and the next image.  Every image is
flat, including the initial copy of $K^0$.  Every such short exact
sequence is tensor-exact.  Therefore (GSL.4) stays exact after any
tensor product, and its lifting characterization again says that
$1\otimes\phi$ is a quasi-isomorphism.

Finally $K^0$ is finite over the Noetherian ring $R$, hence finitely
presented.  A finitely presented flat module is projective: localize
at a prime, lift a basis of its residue vector space, use flatness to
show that the kernel of the resulting finite free surjection has
zero residue module, and apply Nakayama to this finite kernel.
Thus the localization is free.  Finite presentation spreads this
basis to a principal neighborhood.  On a finite such cover, clear
denominators in the local dual-basis identities; powers of the
covering elements still generate the unit ideal.  A linear
combination of these identities is a global dual-basis identity,
which splits a finite free surjection onto the module.
This proves projectivity and applies to $K^0$.  The positive
terms were already finite free.  Thus the exact universal
finite-projective-complex theorem holds, including its
tensor-by-any-module clause.
\end{proof}
% END GSL7

% BEGIN GSL8
\paragraph{GSL8: arbitrary-base transfer from the original Noetherian cube model.}
\label{claim:GSL8}
Let $I$ be a finite linearly ordered set, $R_0\to R$ a ring map
with $R_0$ Noetherian, and $C_0$ and $A$ cubes of modules over
$R_0$ and $R$, respectively.  Suppose a cube map
$\theta:R\otimes_{R_0}C_0\to A$ is an isomorphism at every
nonempty subset of $I$, each corresponding $C_0(J)$ is flat,
and every cohomology module of the unaugmented \v Cech complex
of $C_0$ is finite over $R_0$.  There is a finite-projective
nonnegative complex $K$ over $R$ and a map
$\phi:K\to\check C(A)$ that is a quasi-isomorphism after every
ring map out of $R$.  Moreover $K^p=0$ for $p\ge |I|$.

\begin{proof}
The degree-$p$ term of $\check C(C_0)$ is the finite direct sum of
$C_0(J)$ over the $(p+1)$-element subsets $J$.  It is flat and
vanishes for $p\ge |I|$.  If $I$ is empty, take the zero complex
and zero map.  Otherwise GSL6--GSL7 provide a finite-projective
complex $K_0$ and a universally tensor-valid quasi-isomorphism
$K_0\to\check C(C_0)$, with the same upper bound.  Put
$K=R\otimes_{R_0}K_0$.  Finite projectives remain finite
projective under arbitrary extension of scalars, since their split
finite free presentations tensor to split presentations.

Tensor product commutes with these finite direct sums and their
alternating restriction differentials.  Hence $\theta$ supplies
an actual isomorphism
$R\otimes_{R_0}\check C(C_0)\simeq\check C(A)$.  The empty
vertex is not a term of this unaugmented complex, so no assertion
about $\theta(\varnothing)$ is needed.  Compose the base-changed
map with this isomorphism.  If $R\to R'$ is any subsequent map,
associativity identifies the resulting map with the one obtained by
tensoring $K_0\to\check C(C_0)$ directly by $R'$ over $R_0$.
GSL7 proves it is a quasi-isomorphism.  This proves the source's
full \texttt{exists\_finite\_projective\_complex\_of\_noetherian\_model}
conclusion without a flatness hypothesis on either ring map.
\end{proof}
% END GSL8

% BEGIN GSL9
\paragraph{GSL9: the actual section \v Cech complex under arbitrary base change.}
\label{claim:GSL9}
Let $X\to\operatorname{Spec}R$ have a finite ordered open cover
$D_i$, with every nonempty finite intersection affine, and let $M$
be an actual line bundle.  The complex
\[
 C^p=\bigoplus_{i_0<\cdots<i_p}
       \Gamma(D_{i_0}\cap\cdots\cap D_{i_p},M)
 \tag{GSL.5}
\]
with alternating restriction differential commutes with every ring
base change.  Its degree-zero kernel is the module of actual global
sections.  If $\mathrm V(M)$ holds, it is exact in degree one.
All these identifications respect restrictions and the actual local
coordinates of sections.

\begin{proof}
On an affine intersection $D=\operatorname{Spec}A$, the line bundle
is represented by its invertible $A$-module $Q=\Gamma(D,M)$.
For any $R\to R'$, its pullback to
$\operatorname{Spec}(R'\otimes_R A)$ is represented by
\[
 (R'\otimes_R A)\otimes_AQ\simeq R'\otimes_RQ.
\]
This follows directly by localizing on charts where $Q$ is free of
rank one and gluing the transition units; no flatness of $R'$ is
needed.  These comparisons commute with restriction to smaller affine
intersections, so they give an isomorphism of the whole complexes.

A degree-zero cycle is a family of sections on the $D_i$ agreeing
on pairwise intersections.  GSL1 glues it uniquely to a global
section.  For the degree-one statement, extend an alternating
\v Cech cocycle $\alpha_{ij}$ to all ordered pairs, putting
$\alpha_{ii}=0$ and $\alpha_{ji}=-\alpha_{ij}$.  Its cocycle
equation reads $\alpha_{ij}+\alpha_{jk}=\alpha_{ik}$ as sections
of the same line bundle on a triple intersection.  Refine each $D_i$
by $D_i\cap U_a$, where $U_a$ is an $M$-chart, and express
$\alpha_{ij}$ in the first chart.  This is exactly a cocycle in
the definition of $\mathrm V(M)$.  Its solution gives sections
$\beta_{ia}$ on the refinement.  For fixed $i$, the differences
between $\beta_{ia}$ and $\beta_{ib}$ are $\alpha_{ii}=0$,
so they glue to $\beta_i$ on $D_i$.  Then
$\alpha_{ij}=\beta_i-\beta_j$, proving exactness.  This argument
also explains why the literal all-cover version of $\mathrm V$
is sufficient even when the $D_i$ themselves do not trivialize $M$.
\end{proof}
% END GSL9

% BEGIN GSL10
\paragraph{GSL10: constructing every field of the literal \texttt{Perf} record.}
\label{claim:GSL10}
In GSL9 suppose there is a nonnegative complex $P^\bullet$ of
finite projective $R$-modules and a map
$\phi:P^\bullet\to C^\bullet$ that is a quasi-isomorphism after
every ring map out of $R$.  Then the source's actual
$\texttt{Perf}(f,M)$ record exists.  In particular this conclusion
holds from the Noetherian cube-model data of GSL8 for the actual
section cube of GSL9.

\begin{proof}
Apply GSL5 to $P^0\to P^1\to P^2$, obtaining matrices
$D_0,D_1$ and fixed elements $u_k\in P^0$.  For each cover index
$i$ and each line-bundle chart $j$, use
\[
 W_{ij}=D_i\cap U_j,\qquad \operatorname{ch}(i,j)=j.
\]
These opens cover $X$.  The $i$-component of $\phi^0(u_k)$ is
an actual section of $M$ on $D_i$; let $u_{k,ij}$ be its coordinate
function on $W_{ij}$.  These functions, fixed over the original
ring, are the \texttt{Perf.u} data.

For any $R\to R'$ and any specified Cartesian model
$X'\to X$ of the base change, GSL5 identifies
$\ker(D_0)_{R'}$ with $\ker(d_P^0)_{R'}$.  The assumed
quasi-isomorphism identifies the latter with $\ker d_C^0$
(degree zero has no incoming boundary).  GSL9 identifies that kernel
with $\Gamma(X',M')$, where $M'$ is the actual pullback.  Their
composite is the required set equivalence $E_{R'}$.  On the inverse
image of $W_{ij}$ its coordinate is precisely
\[
 E_{R'}(v)|_{W_{ij}'}
   =\sum_k c_{f'}(v_k)\,\pi^*u_{k,ij}. \tag{GSL.6}
\]
The equality holds because $\phi$ is a linear map, its scalar
extension acts by tensoring, and the GSL9 comparison uses the
actual pulled coordinates.  Thus it is the printed
\texttt{Perf.sec} equality, not merely an abstract dimension
comparison.

If $\mathrm V(M')$ holds, GSL9 gives $H^1(C_{R'})=0$.
The universal quasi-isomorphism gives exactness of
$P^0_{R'}\to P^1_{R'}\to P^2_{R'}$, and GSL5 gives exactness
of the two matrices.  This is exactly \texttt{Perf.exact}.
The fields $a,b,c,D_0,D_1,\mathrm{comp},I,W,\mathrm{ch},
\mathrm{le},\mathrm{cover},u,\mathrm{sec},\mathrm{exact}$ have
all been constructed.  Every assertion holds for arbitrary $R'$ and
for every specified pullback square, since its unique comparison
isomorphism with the standard tensor-product model transports the
same chart equalities.

To apply this result to a \texttt{VeryAmpleBy} presentation,
take $D_i=D(s_i)$.  On $D_i$ the nonvanishing section $s_i$
trivializes the presenting line bundle, so
$D_i\cap D_j$ is the principal open defined by $s_j/s_i$ in
the affine $D_i$.  More generally every nonempty finite
intersection is obtained by inverting a product of these
ratios in one $D_i$, hence is affine.  The finite cover and
all hypotheses of GSL9 therefore hold.  Consequently the
source's universally valid finite-projective model for this
actual section complex, whenever supplied, proves the whole
\texttt{nonempty\_perf\_of\_smooth\_isProjBy} implication.
\end{proof}
% END GSL10

% BEGIN GSL11
\paragraph{GSL11: linearity and naturality forced by \texttt{Perf}'s coordinates.}
\label{claim:GSL11}
Let an actual $\texttt{Perf}(f,M)$ be given, without choosing a
preferred family of its set equivalences.  Its coordinate equation
forces every equivalence $E_{R'}$ to be linear and uniquely determines
it.  For a subsequent ring map $R'\to R''$, the pullback of
$E_{R'}(v)$ is $E_{R''}(v_{R''})$.  The same statements hold after
the literal \texttt{LPow} and \texttt{LinComb} chart conversions.

\begin{proof}
The inverse images of the fixed $W_i$ cover every pullback scheme.
The right side of (GSL.6) is additive in $v$ and linear over the
base ring.  GSL1 therefore identifies the section for $v+v'$ with
the sum of the sections for $v,v'$, and the section for $av$ with
$a$ times the section for $v$.  It also proves that two purported
sections with those coordinates are equal.  Thus the set equivalence
is both linear and uniquely determined by its coordinate formula.

Pull (GSL.6) through the subsequent map.  Pullback commutes with the
finite sum, products, structure maps and the fixed functions $u_{ki}$.
The resulting expression is exactly (GSL.6) for the base-changed
vector.  Uniqueness gives the asserted naturality.  The chart maps
in GSL1 preserve this equation by restrictions, hence preserve its
uniqueness and all these operations.  In particular, the argument
does not assume that an arbitrarily selected equivalence in the
definition of \texttt{Perf} was natural to begin with.
\end{proof}
% END GSL11

% BEGIN GSL12
\paragraph{GSL12: a single local section basis valid on every affine test.}
\label{claim:GSL12}
Let $F=\texttt{Perf}(f,M)$ and $\gamma:R\to K$ be a field point
at which the actual pulled line bundle has $\mathrm V(M_K)$.
There are $r\in R$ with $\gamma(r)\ne0$ and actual sections
$b_1,\ldots,b_m\in\Gamma(X,M)$ whose restrictions form a basis
after \emph{every} map $R\to R'$ inverting $r$.  Here
$m=\dim_K\Gamma(X_K,M_K)$.  If $R$ is local and $\gamma$ factors
through an extension of its residue field, they already form a basis
over $R$.

\begin{proof}
The \texttt{Perf.exact} field gives middle exactness at $K$ for
$D_0,D_1$.  Apply GSL3--GSL4, and let $b_i$ be the actual section
represented by the $i$th column of $B$.  This column is a cycle
over $R$ itself, since $D_0B=0$ before localization.  GSL11 says
that any actual pulled section satisfying \texttt{IsPullSec} or
\texttt{IsPullSecLC} is represented by the pulled column.  GSL4
gives unique coordinates in these columns for every kernel vector;
GSL11 transfers them, with their uniqueness, to the actual section
module.  Thus the sections are independent and span on every test.
At $K$ they form a basis of length $m$, so that length is the stated
dimension.  The final assertion is the local-unit clause of GSL4.
\end{proof}
% END GSL12

% BEGIN GSL13
\paragraph{GSL13: lifting actual sections through a surjective coefficient map.}
\label{claim:GSL13}
Let $F=\texttt{Perf}(f,M)$, where $M$ may be the literal
$\operatorname{LinComb}(L,\phi,\epsilon)$, and let
$\rho:R\to\overline R$ be a surjective ring map.  Suppose a fixed
matrix $B$ has $D_0B=0$ over $R$ and its columns span
$\ker(D_0)_{\overline R}$.  Then every section of the actual pulled
$M$ over $\overline R$ lifts to a section over $R$.  For a
\texttt{LinComb}, the lift satisfies the printed
\texttt{IsPullSecLC} coordinate equalities.

\begin{proof}
Use the coordinate equivalence for the specified pullback square
to represent a given section $\overline s$ by
$\overline v\in\ker(D_0)_{\overline R}$.  Write
$\overline v=\overline B\,\overline z$ and lift each coordinate
of $\overline z$ to $z\in R^m$.  The vector $v=Bz$ is a cycle
over $R$.  The identity-base-change equivalence, followed by the
inverse of the harmless identity-pullback chart conversion in GSL1,
gives an actual section $s$ of $M$.  By GSL11 its pullback has the
same coordinates as $\overline s$ on the inverse images of all
$W_i$, hence equals $\overline s$.

For clarity, if $M=\operatorname{LinComb}(L,\phi,\epsilon)$ and
$\pi$ is the map on total spaces, the target in the source is the
line bundle with maps $\pi\phi_k$ on the same finite index set, not
the syntactically nested bundle $M.\mathrm{pull}(\pi)$.  The chart
conversions of GSL1 identify these bundles by the identity map on
each common coordinate function.  Accordingly the equality just
proved reads, for every such chart $j$,
\[
 \overline s_j
   =\operatorname{res}\bigl(\pi^*s_j\bigr).
 \tag{GSL.7}
\]
This is exactly \texttt{IsPullSecLC}.  Surjectivity of $\rho$ and
the single fixed spanning matrix are all that the argument needs;
neither locality nor a field hypothesis enters this lifting lemma.
\end{proof}
% END GSL13

\subsection{The geometric inputs and their proved reductions}

The next arguments use the original \texttt{PolAV} record $P/S$.
It contains a commutative abelian scheme $A/S$, its relative dimension
$g$, and an actual rigidified symmetric relatively ample line bundle
$L$.  Its \texttt{KTriv} condition tests the translation kernel over
every affine base: local triviality of $t_x^*L\otimes L^{-1}$ forces
$x$ to be the identity.  In particular we do not replace this record
by the data of an arbitrary ample bundle or by complex fibers alone.

We state the upstream inputs when they are used.  The original
\texttt{PolBaseChangeSt}, denoted $\mathcal P$, says: for every map
$t:S'\to S$ and every $P/S$, there is a polarized record $P'/S'$
and a Cartesian map $G:A'\to A$ with $g'=g$, preserving multiplication
on every test scheme, together with an actual generated
line-bundle isomorphism $G^*L\simeq L'$.  This entire record remains
an input here.

% BEGIN GSL14
\paragraph{GSL14: transport of sections, bases and the all-cover vanishing law.}
\label{claim:GSL14}
Let $e:X'\to X$ be a scheme isomorphism and let $M',M$ be actual
line bundles, with a generated isomorphism $M'\simeq e^*M$.
It induces a linear equivalence of section modules, preserving
specified coefficient-ring actions when $e$ is over the base.
It transports bases and their unique coordinate expressions.
Moreover $\mathrm V(M)$ implies $\mathrm V(M')$, with the exact
all-cover meaning used above.

\begin{proof}
Refine the charts of $M'$ by inverse images of charts of $M$.
On a common chart the generated isomorphism is multiplication by
a specified unit $u$.  Send a coordinate function $s$ to
$u\,e^*s$ (or to $u^{-1}e^*s$, according to the displayed direction
of the isomorphism).  The defining unit compatibility is exactly
the equation ensuring that these functions form a section.
GSL1 glues them.  The inverse scheme map and inverse units give
the inverse section map.  Both maps are linear, commute with the
base structure maps, and are mutual inverses on every chart.
They therefore transport bases and unique coefficients.

For vanishing, start with a cocycle for $M'$ on an arbitrary
subordinate cover $V_a$.  Refine it by the inverse images of
the $M$-charts and transport the cocycle through the inverse units.
The unit compatibility turns its cocycle equation into the one
for $e^*M$.  Transporting across the isomorphism $e$ gives a
cocycle for $M$; solve it using $\mathrm V(M)$ and pull the
solution back.  Multiply by the original units.  This gives a
solution on the refinement for the initial cocycle.

For fixed $a$, its solutions on two refined pieces of $V_a$
agree: their required difference is the diagonal cocycle
$\phi_{aa}$, which is zero by the cocycle equation with three
equal indices and $g_{i(a)i(a)}=1$.  Thus the solutions glue to
one function on each $V_a$.  The coboundary equation holds
locally and hence on every $V_a\cap V_b$.  This proves the
original \texttt{h1Zero\_of\_isIso} statement, including changes
of both the scheme and the literal line-bundle charts.
\end{proof}
% END GSL14

% BEGIN GSL15
\paragraph{GSL15: the required projective presentation after any affine base change.}
\label{claim:GSL15}
Assume $\mathcal P$.  For every $P/S$, every commutative $R$,
$t:\operatorname{Spec}R\to S$, and every specified Cartesian model
\[
 \begin{array}{ccc}
 X&\xrightarrow{\pi}&A\\
 f\downarrow&&\downarrow\\
 \operatorname{Spec}R&\xrightarrow{t}&S,
 \end{array}
\]
there exist $N\ge1$, $m\ge0$ and $m+1$ actual sections of
$(\pi^*L)^N$ giving the source's \texttt{VeryAmpleBy} presentation
of $f$.  In particular $f$ satisfies \texttt{IsProjBy}.

\begin{proof}
Apply $\mathcal P$ to obtain $P'/\operatorname{Spec}R$.
The two Cartesian models are uniquely isomorphic over
$\operatorname{Spec}R$, and the generated isomorphism in
$\mathcal P$ identifies their actual pulled line bundles.
Apply the \texttt{RelAmple} field of $P'$ on the whole affine
base $\operatorname{Spec}R$.  It supplies a positive power,
$m+1$ sections whose nonzero loci cover, an affine nonzero
locus for each section, and on each such locus ratios $r_j$
with $s_j=r_js_i$ and a surjection
\[
 R[X_0,\ldots,X_m]\longrightarrow
 \Gamma(D(s_i),\mathcal O),\qquad X_j\longmapsto r_j.
 \tag{GSL.8}
\]
This is the original definition of relative ampleness used here.

Transport all these sections to $X$ by GSL14, also converting
the repeated-factor power to the actual $\operatorname{pw}$
charts.  Multiplication by a transition unit preserves whether
the germ of a section is a unit; consequently the nonzero
loci are the inverse images of the old ones.  They cover and
remain affine under the scheme isomorphism.  The section
isomorphism multiplies numerator and denominator by the
same unit, so the ratios are pulled functions $r_j$.
The corresponding map (GSL.8) is conjugate through ring
isomorphisms to the old surjection.  It remains surjective.
These are every field of \texttt{VeryAmpleBy}, with the same
positive integer $N$.  No restriction on $R$ or $S$ was imposed.
\end{proof}
% END GSL15

% BEGIN GSL16
\paragraph{GSL16: descent of first-cocycle vanishing by a finite average.}
\label{claim:GSL16}
Let $f:Y\to X$, let a finite group $G$ act on $Y$ by maps
$\tau_g$ over $X$, with the composition convention
$\tau_g\tau_h=\tau_{gh}$, and suppose $|G|$ is a unit in
$\Gamma(X,\mathcal O)$.  Assume that on every open $U\subseteq X$
pullback identifies $\Gamma(U,\mathcal O)$ with the functions on
$f^{-1}U$ invariant under every $\tau_g$: each invariant function
has a unique preimage.  If $\mathrm V(f^*M)$ holds, then
$\mathrm V(M)$ holds.

\begin{proof}
First this function-descent premise makes every pullback
$\Gamma(U,\mathcal O)\to\Gamma(f^{-1}U,\mathcal O)$ injective.
Indeed a pulled function is invariant because $f\tau_g=f$,
and uniqueness of its descended preimage proves injectivity.

Pull an arbitrary cocycle $\phi_{ab}$ for $M$ on a subordinate
cover $V_a$ to $Y$.  Vanishing there gives
\[
 f^*\phi_{ab}=\psi_a-f^*g_{i(a)i(b)}\psi_b.
\]
Set $\Psi_a=\sum_{g\in G}\tau_g^*\psi_a$.  Reindexing the sum
by group multiplication proves that $\Psi_a$ is invariant,
so it descends uniquely to a function $\chi_a$ on $V_a$.
Every pulled transition function and every $f^*\phi_{ab}$ is
invariant.  Sum the displayed coboundary identity, and use the
injectivity just proved to descend the resulting equation:
\[
 |G|\phi_{ab}=\chi_a-g_{i(a)i(b)}\chi_b.
\]
Multiply each $\chi_a$ by the restriction of $|G|^{-1}$.
The resulting functions solve the original cocycle.  The argument
uses every original cover without a refinement or a separability
assumption on $f$ beyond the stated invariant-function law.
\end{proof}
% END GSL16

% BEGIN GSL17
\paragraph{GSL17: positive-power vanishing from the original geometric suppliers.}
\label{claim:GSL17}
Assume $\mathcal P$ and the following original inputs.

\begin{enumerate}
 \item \emph{Serre input.}  For every Noetherian $R$, every actual
 $\texttt{VeryAmpleBy}(f,H,s)$ and every actual line bundle $M$,
 there is $t_0\in\mathbb Z$ such that
 $\mathrm V(M\otimes H^t)$ for all integers $t\ge t_0$.
 \item \emph{Multiplication pullback.}  For every $S,P,n$ there
 is an actual generated trivialization of
 $L^{n^2}\otimes([n]^*L)^{-1}$.
 \item \emph{Multiplication geometry.}  For every locally Noetherian
 $S$, $P/S$ and $n\ge1$, the map $[n]$ is finite, flat and
 surjective, and is etale if $n$ is a unit on $S$.
 \item \emph{Torsion-invariant descent.}  For every algebraically
 closed $K$, $P/K$ and $n$ nonzero in $K$, functions on
 $[n]^{-1}U$ invariant under all translations by $A[n](K)$
 descend uniquely to $U$, for every open $U\subseteq A$.
\end{enumerate}
Then for every algebraically closed field $K$, every original
$P/K$ and every $n\ge1$, one has $\mathrm V(L^n)$.

\begin{proof}
By GSL15 choose $N\ge1$ with $H=L^N$ very ample in the printed
sense.  Apply the Serre input to $L^r$ for each
$0\le r<N$.  Choose a nonnegative integer $T$ at least all
the resulting bounds.  Choose a prime $q$ sufficiently large
that $q$ is nonzero in $K$ and $q^2n/N\ge T+1$.
There are arbitrarily large primes, so this choice works also
in positive characteristic by avoiding that characteristic.
Write
\[
 q^2n=r+Nt,\qquad 0\le r<N.
\]
Then $t\ge T$ and the Serre input gives
$\mathrm V(L^r\otimes(L^N)^t)=\mathrm V(L^{q^2n})$.
The multiplication-pullback generator, raised to its $n$th
power and combined with the identity $q^2n=r+Nt$, gives an
actual generated isomorphism
$[q]^*L^n\simeq L^{q^2n}$.  On common charts its unit is the
$n$th power of that generator times the necessary transition
units; equality of the exponents makes the cocycle equation
identical on both sides.  GSL14 therefore gives
$\mathrm V([q]^*L^n)$.

Let $G=A[q](K)$.  It is finite: its points belong to the finite
fiber over the identity of the finite map $[q]$.  It is an
abelian group killed by the prime $q$, hence a finite vector
space over $\mathbb F_q$ and has order $q^a$ for some $a\ge0$.
Because $q\ne0$ in $K$, this order is a unit in $K$ and hence
in $\Gamma(A,\mathcal O)$.  Translation by $G$ is an action
over $[q]$, since $[q](x+y)=[q]x$ for $y\in G$.
The invariant-function input is exactly the descent premise
of GSL16.  Apply that claim with $f=[q]$ and $M=L^n$.
It yields $\mathrm V(L^n)$.  This proves the entire original
positive-power vanishing supplier, uniformly over all
algebraically closed characteristics, from the four displayed
inputs; no characteristic-zero vanishing theorem was inserted.
\end{proof}
% END GSL17

% BEGIN GSL18
\paragraph{GSL18: the exact exponent forced by growth and square dilation.}
\label{claim:GSL18}
Let $h:\mathbb N\to\mathbb N$ be nondecreasing, $h(1)=1$, and
$G\subseteq\mathbb N$ satisfy $0\notin G$ and, for every $m$,
at least one of $m,m+1$ belongs to $G$.  Suppose
\[
 h(n^2j)=h(n^2)h(j)\quad(n\in G,\ j\ge1), \tag{GSL.9}
\]
and there are $N\ge1$, $c,C\in\mathbb N$ and $g\ge0$ such that
\[
 t^g\le c\,h(Nt),\qquad h(Nt)\le Ct^g\quad(t\ge1).
 \tag{GSL.10}
\]
Then $h(j)=j^g$ for every $j\ge1$.

\begin{proof}
Fix $n\in G$ and put $a=n^2\ge1$.  Iterating (GSL.9) gives
$h(a^ki)=h(a)^kh(i)$ for every $i\ge1$.  By monotonicity and
the upper bound in (GSL.10),
\[
 h(a)^k=h(a^k)\le h(Na^k)\le C(a^g)^k.
\]
Since a ratio greater than one has unbounded powers, this
forces $h(a)\le a^g$.  The lower bound in (GSL.10) gives
\[
 (a^g)^k\le c\,h(Na^k)=c\,h(N)h(a)^k.
\]
Here $h(a)\ge1$, so the same argument forces $a^g\le h(a)$.
Thus $h(n^2)=n^{2g}$ for every $n\in G$.

There are arbitrarily large $n\in G$.  Fix $j\ge1$ and for each
such $n$ choose $m_-,m_+\in G$ immediately below and above
$n\sqrt j$.  They can be chosen with
$0\le n\sqrt j-m_-\le2$ and
$0\le m_+-n\sqrt j\le2$, because $G$ meets every pair of
consecutive integers.  For large $n$ both are positive.
Monotonicity and (GSL.9) yield
\[
 m_-^{2g}\le n^{2g}h(j)\le m_+^{2g}.
\]
Divide by $n^{2g}$ and let these $n$ tend to infinity.
Both bounds tend to $(\sqrt j)^{2g}=j^g$.
This proves the assertion, including $g=0$.  The use of real
limits merely shortens the source's equivalent integer estimates.
\end{proof}
% END GSL18

% BEGIN GSL19
\paragraph{GSL19: the original geometric rank is exactly \(n^g\).}
\label{claim:GSL19}
Assume the following source inputs for every original $P/K$ over
an algebraically closed field.

\begin{enumerate}
 \item Its geometric total space is integral (the source obtains this
 from smoothness and geometric connectedness).
 \item $h^0(L)=1$.
 \item For $n\ne0$ in $K$ and $j\ge1$ there is an actual
 coefficient-linear equivalence
 \[
 \Gamma(A,L^{n^2j})\simeq
       \prod_{x\in A[n](K)}\Gamma(A,L^j).
 \tag{GSL.11}
 \]
 \item Every $\Gamma(A,L^j)$ for $j\ge0$ is finite, and for some
 $N,c,C\in\mathbb N$ with $N,c\ge1$ its dimensions obey
 (GSL.10) with the original $g=P.g$.
\end{enumerate}
Then \texttt{GlobalSectionsRankAlgClosed} holds: for every such
$P$ and every $j\ge1$ the literal \texttt{SecM} has dimension $j^g$
over $\Gamma(\operatorname{Spec}K,\mathcal O)$.

\begin{proof}
Use GSL1 and the identity-base product isomorphism to identify
the displayed modules with the literal \texttt{SecM} modules,
preserving the canonical base coefficient ring.  Put
$h(j)=\dim_K\Gamma(A,L^j)$.  The rank-one input gives a nonzero
actual section $s$ of $L$.  Multiplication by $s$ maps
$\Gamma(A,L^j)$ injectively into $\Gamma(A,L^{j+1})$.
Indeed, on an integral scheme a nonzero section of a line bundle
is nonzero in every nonempty trivializing open: otherwise its
zero coordinate, transferred across a nonempty intersection,
would make it zero in every chart.  The coordinate rings of
nonempty opens of an integral scheme are domains, so
$fs=0$ forces every local coordinate of $f$ to vanish.
Thus $h$ is nondecreasing, using the finite-dimensionality input.

Let $G=\{n:(n:K)\ne0\}$.  It omits zero and meets each pair
of consecutive integers, since both cannot vanish in a field.
For $n\in G$, the $j=1$ instance of (GSL.11) implies that
$A[n](K)$ is finite: otherwise its product with the nonzero
one-dimensional space $\Gamma(A,L)$ would contain arbitrarily
large independent families of coordinate vectors, contradicting
the finite dimension of $\Gamma(A,L^{n^2})$.
The same instance gives $h(n^2)=|A[n](K)|$.  The general $j$
instance therefore gives
$h(n^2j)=|A[n](K)|h(j)=h(n^2)h(j)$.
Apply GSL18 with the stated growth bounds.  The resulting
$h(j)=j^g$ is preserved by the canonical coefficient-ring
isomorphism, proving the printed rank law.  The proof uses the
actual section equivalences, not an independently assumed
degree formula for $[n]$.
\end{proof}
% END GSL19

% BEGIN GSL20
\paragraph{GSL20: rank and vanishing on every actual geometric fiber.}
\label{claim:GSL20}
Assume $\mathcal P$, positive-power vanishing on every $P'/K$,
and \texttt{GlobalSectionsRankAlgClosed}.  Then for every scheme
$S$, original $P/S$, algebraically closed $K$, point
$t:\operatorname{Spec}K\to S$ and $n\ge1$,
\[
 \dim_{\Gamma(\operatorname{Spec}K,\mathcal O)}
    \operatorname{SecM}(P.A,P.L,n,\operatorname{Over.mk}t)
       =n^{P.g}. \tag{GSL.12}
\]
The actual positive power on any specified Cartesian model of this
fiber has $\mathrm V$.  The same assertions apply after any
intermediate base scheme $T\to S$ and any geometric point of $T$.

\begin{proof}
Apply $\mathcal P$ to $t$ and obtain $P'/K$, its Cartesian
map $G:A'\to A$, $P'.g=P.g$, and its actual line-bundle
isomorphism.  Every specified fiber model is uniquely
isomorphic over $K$ to $A'$.  Raise the generated isomorphism
to the $n$th power and apply GSL14.  This transports both
the section module and the all-cover vanishing assertion
from $(L')^n$.  The one-factor \texttt{LinComb} chart maps
of GSL1 identify the section module with the exact
\texttt{SecM} in (GSL.12), so the closed-field rank law gives
the stated integer.

For $T\to S$ and $\sigma:\operatorname{Spec}K\to T$, apply
the same construction to their composite.  The universal
property of the two successive pullbacks gives a Cartesian
map from $A'$ to $T\times_S A$.  Its composite to $A$ is
$G$, so the actual pulled powers and their chart isomorphisms
are precisely those already used.  Neither $S$ nor $T$ has
been assumed Noetherian or reduced.  This proves the exact
full-fiber rank and vanishing inputs of the two section-law
theorems, not only their instances at a chosen affine chart.
\end{proof}
% END GSL20

% BEGIN GSL21
\paragraph{GSL21: the geometric root and the actual translated-power generator.}
\label{claim:GSL21}
Assume the multiplication geometry of GSL17 and the original
square input: for every $S,P$, every $T\to S$, points $x,y\in A(T)$
and point of $T$, the bundle
\[
 t_{x+y}^*L\otimes L\otimes(t_x^*L)^{-1}
                    \otimes(t_y^*L)^{-1}
 \tag{GSL.13}
\]
has a generated trivialization after restriction above some open
neighborhood of that point.  Here the notation is on $T\times_S A$.
Then, for every locally Noetherian $S$, every $P/S$, every
algebraically closed field $K$, $t:\operatorname{Spec}K\to S$,
$q\ge1$ and $c\in A(K)$ above $t$, there exist $d$ with $qd=c$
and an actual generated trivialization of
\[
 L^{q-1}\otimes t_c^*L\otimes(t_d^*L)^{-q}. \tag{GSL.14}
\]
This is the full original \texttt{twisted\_pow\_translate} supplier.

\begin{proof}
The fiber of the finite surjective map $[q]$ above the $K$-point
$c$ is a nonempty finite $K$-scheme.  It has a closed point with
residue field finite over $K$, hence equal to $K$ because $K$ is
algebraically closed.  Its point gives $d$ with $qd=c$.  This
argument does not require $q$ to be invertible in $K$.

On $\operatorname{Spec}K$ any neighborhood of its unique point
is the whole base.  Thus the square input gives an actual
global generator for (GSL.13) for each pair of $K$-points.
Put $M_x=t_x^*L$ on the geometric fiber.  The generator
identifies $M_{x+y}\otimes M_0$ with $M_x\otimes M_y$.
Multiplying these generators successively with $(x,y)=(ad,d)$
proves by induction that
\[
 M_{qd}\otimes M_0^{q-1}\otimes M_d^{-q}
\]
has an actual generator.  The starting case $q=1$ is the
identity generator; the step multiplies the previous generator
by the square generator and cancels the common factors.
All cancellations are products or inverses of the specified
units on a common chart, so they yield a generated
\texttt{LinComb} isomorphism, not just a numerical identity.
Since $qd=c$, this is (GSL.14) with the exact maps
$[\mathrm{snd},\mathrm{translPt}(c),\mathrm{translPt}(d)]$
and exponents $[q-1,1,-q]$.
\end{proof}
% END GSL21

% BEGIN GSL22
\paragraph{GSL22: a pure-power model of each twisted geometric fiber.}
\label{claim:GSL22}
Assume $\mathcal P$ and the conclusion of GSL21 at its printed
locally Noetherian scope.  Let $S$ now be \emph{arbitrary},
$P/S$ original, $T\to S$ any base, $c\in A(T)$, $q\ge1$ and
$\sigma:\operatorname{Spec}K\to T$ an algebraically closed point.
Put $X=T\times_S A$ and
$M=L_T^{q-1}\otimes t_c^*L$ on $X$.  There are a polarized
$P'/K$ with $g'=g$ and a Cartesian map
$\pi:A'\to X$ over $\sigma$ for which
\[
 \pi^*M\simeq (L')^q \tag{GSL.15}
\]
is an actual generated isomorphism.  One may equally use the
literal $\operatorname{baseOver}(K)\times_K A'$ model and its
$\operatorname{LPow}$ in (GSL.15).  Thus, under the rank and
vanishing inputs of GSL20, $\pi^*M$ has dimension $q^g$ and
has $\mathrm V$.

\begin{proof}
Apply $\mathcal P$ to the composite $\operatorname{Spec}K\to S$,
obtaining $P'/K$, its multiplication-preserving Cartesian
map $G:A'\to A$, and $G^*L\simeq L'$.  The point $c$ lifts
uniquely to $c'\in A'(K)$.  The resulting Cartesian map
$j:A'\to X$ satisfies, on every test scheme,
\[
 j^*L_T\simeq L',\qquad
 j^*(t_c^*L)\simeq t_{c'}^*L'.
 \tag{GSL.16}
\]
The second equality uses the actual multiplication-preserving
field of $\mathcal P$: $G(x+c')=G(x)+c$.
Take the powers and tensor products of its given units to
obtain these actual generated isomorphisms.

Apply GSL21 to the new base $\operatorname{Spec}K$ and $P'$.
This base is Noetherian regardless of the original $S$.
It gives $d\in A'(K)$ with $qd=c'$ and a generated
isomorphism
$(L')^{q-1}\otimes t_{c'}^*L'\simeq t_d^*((L')^q)$.
Let $\tau=t_{-d}:A'\to A'$.  It is an isomorphism over $K$
with inverse $t_d$.  Set $\pi=j\circ\tau$.
Precomposing a Cartesian square by this isomorphism over the
base keeps it Cartesian.  Pulling the preceding line-bundle
isomorphism by $\tau$ gives (GSL.15), because
$t_d\circ t_{-d}=1$.

For the source's literal \texttt{baseOver} product, compose
with its canonical isomorphism to $A'$ and use the one-factor
chart maps of GSL1.  This is the exact source construction
$\pi=\tau\gg j$, including its inverse-translation sign.
GSL14 transports the rank and all-cover vanishing of $(L')^q$
to this actual $\pi^*M$.  In particular no locally Noetherian
assumption on $S$ has entered through the geometric twist input.
\end{proof}
% END GSL22

\subsection{The original global laws, with their exact quantifier order}

The remaining geometric model premise, denoted $\mathcal F$, is the
literal source statement
\[
 \begin{gathered}
 \text{for every commutative }R,\quad
 f:X\to\operatorname{Spec}R\text{ smooth and }\texttt{IsProjBy},\\
 \text{and every actual line bundle }M,\quad
        \operatorname{Nonempty}(\texttt{Perf}(f,M)).
 \end{gathered} \tag{GSL.17}
\]
GSL5--GSL10 prove its complete finite-complex construction once the
actual section cube has the stated Noetherian flat model with finite
cohomology.  The global existence of that geometric model is a retained
supplier.  More precisely, the original source supplies it through
\texttt{exists\_noetherian\_model\_grCube} and the dehomogenization
comparison with the actual section cube, followed by
\texttt{exists\_finite\_projective\_complex\_cechLB}.  Its graded
presentation, smooth finite-presentation descent and finite graded
cohomology inputs are not proved by the matrix argument alone.

Write $\mathcal V$ for positive-power $\mathrm V$ over every
algebraically closed field, $\mathcal R$ for
\texttt{GlobalSections\allowbreak RankAlgClosed}, and $\mathcal T$ for the
full locally Noetherian geometric translated-power conclusion in
GSL21.  GSL17, GSL19 and GSL21 respectively give their complete
reductions from the displayed original geometric inputs.

% BEGIN GSL23
\paragraph{GSL23: both original section base-change laws.}
\label{claim:GSL23}
Assume $\mathcal F,\mathcal P,\mathcal V,\mathcal R$.
Then \texttt{SecBaseChangeAnySt} and \texttt{SecBaseChangeSt} both
hold.  Explicitly, for every scheme $S$, every original $P/S$,
$n\ge1$ and point $p\in S$, there exist an open $V\ni p$
and actual sections $b_i$, $i\in\operatorname{Fin}(n^g)$,
of the specified $\operatorname{LPow}$ over $V$, such that
for every commutative ring $R'$ and every
$t:\operatorname{Spec}R'\to V$, \emph{every} tuple $b_i'$
satisfying the literal \texttt{IsPullSec} equalities is
linearly independent and spans the literal \texttt{SecM}
over $\Gamma(\operatorname{Spec}R',\mathcal O)$.

\begin{proof}
Choose an affine open neighborhood of $p$, written
$t_0:\operatorname{Spec}R\hookrightarrow S$, and let
$\mathfrak p\subset R$ be its point above $p$.
On the actual pullback $X=\operatorname{Spec}R\times_S A$
let $M=\operatorname{LPow}(A,L,n,\operatorname{Over.mk}t_0)$.
The morphism $X\to\operatorname{Spec}R$ is smooth by base
change.  GSL15 supplies its \texttt{IsProjBy} presentation.
Thus $\mathcal F$ provides an actual $\texttt{Perf}(f,M)$.

Take $K$ to be an algebraic closure of the residue field
$\kappa(\mathfrak p)$ and let $\gamma:R\to K$ be the
canonical map.  Its kernel is $\mathfrak p$.
GSL20 gives $\mathrm V$ on an actual Cartesian model of
the geometric pullback and gives its section dimension
$n^g$.  Apply GSL12 to obtain a scalar $r\notin\mathfrak p$
and exactly $n^g$ actual sections $b_i^0$ over $X$.
The columns representing these sections were chosen over
$R$ and satisfy $D_0B=0$ there.

Let $V$ be the image of $D(r)\subseteq\operatorname{Spec}R$
under the chosen open immersion.  It contains $p$ and
is canonically isomorphic to $\operatorname{Spec}R[r^{-1}]$.
Pull the $b_i^0$ to the specified restriction object over
$V$ to obtain $b_i$.  For any later $t:\operatorname{Spec}R'\to V$,
the corresponding map $R\to R'$ sends $r$ to a unit.
GSL12 therefore gives unique section coordinates in the
pulled columns.  A tuple $b_i'$ satisfying the original
\texttt{IsPullSec} conditions consists of exactly those
sections by GSL1 and GSL11.  Its unique-coordinate property
is equivalent to linear independence and full span in
$\operatorname{SecM}$: the map from $R'^{n^g}$ taking
coefficients to the linear combination of $b_i'$ is a
linear bijection.  The canonical $\Gamma\operatorname{Spec}$
ring isomorphism gives the printed coefficient action.

The same $V$ and tuple $b_i$ were selected before all
$R',t,b'$.  This proves \texttt{SecBaseChangeAnySt} with
arbitrary original $S$ and arbitrary coefficient ring $R'$.
Restricting the original $S$ to locally Noetherian schemes
gives \texttt{SecBaseChangeSt}.  Thus every premise in the
original \texttt{secBaseChange\_of\_perf} theorem has either
been retained exactly or supplied above: projectivity by
GSL15, isomorphism invariance by GSL14, the universal matrix
theorem by GSL3--GSL4 and the full geometric-fiber rank
statement by GSL20.
\end{proof}
% END GSL23

% BEGIN GSL24
\paragraph{GSL24: the actual twisted basis over every local coefficient ring.}
\label{claim:GSL24}
Assume $\mathcal F,\mathcal P,\mathcal V,\mathcal R,\mathcal T$.
For every arbitrary scheme $S$, original $P/S$, $q\ge1$,
commutative local ring $R$, $t:\operatorname{Spec}R\to S$
and $c\in A(R)$ above $t$, the actual bundle
\[
 M=\operatorname{twistLB}(A,L,q,c)
     =L_R^{q-1}\otimes t_c^*L
\]
has a basis of $q^g$ actual sections $b_i$ over $R$.
For every actual section $s$ there is a unique
$a:\operatorname{Fin}(q^g)\to\Gamma(\operatorname{Spec}R,\mathcal O)$
such that on each literal common chart $j$,
\[
 s_j=\sum_i c_f(a_i)b_{i,j}. \tag{GSL.18}
\]
The same sections remain a basis after every ring map $R\to R'$.

\begin{proof}
Set $X=\operatorname{Spec}R\times_S A$ with its actual projection
$f$.  Smoothness is stable under base change, and GSL15 makes
$f$ \texttt{IsProjBy}.  The tensor formula for $M$ uses the
given transition units, so $M$ is an actual line bundle.
Obtain $F=\texttt{Perf}(f,M)$ from $\mathcal F$.

Let $\kappa$ be the residue field of $R$ and $K$ an algebraic
closure of $\kappa$, with $\gamma:R\to\kappa\hookrightarrow K$.
GSL22 produces a specified Cartesian model of $M_K$ as a
positive power of an original $P'/K$, with dimension $q^g$
and $\mathrm V$.  Apply GSL3--GSL4 to the matrices of $F$
at this geometric model.  The resulting $r$ has nonzero
image in $K$, hence is a unit of $R$, since $\gamma$ factors
through the residue field.  Let $b_i$ be the sections
represented by the fixed columns of $B$.

Over $K$ these sections form a basis.  Their number is
therefore $q^g$ by GSL22, and we reindex the fixed columns
once by $\operatorname{Fin}(q^g)$.  GSL11 and GSL4 give
unique coordinates over $R$, yielding (GSL.18).
Because $r$ is already a unit, its image is a unit under
every ring map $R\to R'$, so the same proof gives the stated
universal basis.  The source's identity-pullback comparison
is removed by the coordinate-preserving chart isomorphism
of GSL1; the resulting $b_i$ belong to the literal
\texttt{twistLB}, and (GSL.18) is its exact \texttt{IsSecBasis}.
\end{proof}
% END GSL24

% BEGIN GSL25
\paragraph{GSL25: every residue section lifts, giving the full \texttt{SecTwistSt}.}
\label{claim:GSL25}
Under GSL24, let $\overline X=\operatorname{Spec}\kappa\times_S A$
and $\rho_X:\overline X\to X$ be the prescribed map induced by
the residue homomorphism $R\to\kappa$.  For \emph{every} actual
section $\overline s$ of the literal bundle
\[
 \operatorname{LinComb}\!\left(
   L,\ [\,\mathrm{snd}\circ\rho_X,
              \mathrm{translPt}(c)\circ\rho_X\,],\
          [\,q-1,1\,]\right),
 \tag{GSL.19}
\]
where composition has its usual right-to-left meaning,
there is a section $s$ of
$\operatorname{twistLB}(A,L,q,c)$ with
$\operatorname{IsPullSecLC}(s,\overline s)$.
Together with GSL24 this is the complete original
\texttt{SecTwistSt}, with no condition on $S$ and no
Noetherian, reduced or flatness condition on the local $R$.

\begin{proof}
The columns $B$ constructed in GSL24 form a basis after the
residue map, because their denominator was a unit of $R$.
The residue homomorphism is surjective.  Apply GSL13 to
this same matrix and the actual bundle $M$.  It lifts every
$\overline s$, with exactly the chart formula (GSL.7).
Equivalently, pull the basis $b_i$ to obtain $\overline b_i$,
write uniquely $\overline s=\sum_i\overline a_i\overline b_i$,
choose arbitrary lifts $a_i\in R$, and set $s=\sum_i a_i b_i$.
Pullback preserves this sum and its coefficients, giving
$\overline s$ on every chart.

The target bundle (GSL.19) is precisely the original
\texttt{LinComb} obtained by composing the two maps with the
whiskered residue map.  GSL1 identifies it with the nested
pullback of $M$ by equality of the local coordinate functions,
so the equality just established is the required
\texttt{IsPullSecLC}, not an equality only after an unspecified
line-bundle isomorphism.  No lift of the geometric root $d$
to $R$ is needed: that root was used only to prove geometric
exactness and rank before choosing the matrix $B$.
The original section, point $c$, local ring and residue
section remain unchanged throughout.
\end{proof}
% END GSL25

% BEGIN GSL26
\paragraph{GSL26: the exact conditional installation and the remaining boundary.}
\label{claim:GSL26}
The literal inputs $\mathcal F,\mathcal P,\mathcal V,\mathcal R$
supply both global section base-change laws; adding the original
$\mathcal T$ supplies the full twisted law.  The constructions
above establish this implication on the original arbitrary bases,
all affine target rings and all local coefficient rings, with
the quantified choices and actual section objects of the source.
Equivalently, $\mathcal V,\mathcal R,\mathcal T$ may be supplied
by the exact upstream packages in GSL17, GSL19 and GSL21.
The finite-complex part of $\mathcal F$ is supplied by
GSL5--GSL10 from its original geometric section-cube model.

Consequently these law instances may be installed in the already
proved TDR/NFR/NFL arguments for the \emph{same} supplied
\texttt{Moduli} record and actual normalized family.
This removes no other premise of those arguments and does not
assert existence of that moduli record or of a new common bundle.

\begin{proof}
GSL23 proves each quantifier of \texttt{SecBaseChangeAnySt};
its restriction proves \texttt{SecBaseChangeSt}.  GSL24 and
GSL25 prove the two conjuncts of \texttt{SecTwistSt}, with
the same $P,q,R,t,c$ in both.  In GSL23 the open and its
basis precede every affine test map.  In GSL24 the geometric
root is allowed only inside the proof, whereas the resulting
sections lie on the original bundle over $R$; GSL25 lifts
every residue section with that same fixed basis.  These
are precisely the clauses consumed by TDR13 and NFL22.
Substitution of the supplied laws in those proved implications
therefore keeps their original objects and quantifier order.

For the upstream form, GSL15 derives the projectivity premise
from $\mathcal P$, GSL14 supplies the full isomorphism-invariance
premise, and GSL3--GSL4 supply the original universally quantified
matrix premise.  GSL17 supplies $\mathcal V$; GSL19 supplies
$\mathcal R$; GSL20 makes their full original-fiber versions.
GSL21 supplies $\mathcal T$, and GSL22 explains why its
locally Noetherian original scope suffices for an arbitrary
$S$ in the final twisted law.  Finally GSL8--GSL10 give
every field of $\mathcal F$ whenever the geometric
Noetherian model of the actual section cube is provided.

The remaining global suppliers are thus explicitly geometric:
the universal section-cube model, polarization base change,
the Serre and multiplication/descent inputs of GSL17,
the initial rank-one, section-decomposition and growth inputs
of GSL19, and the square input of GSL21.  Their formal proofs
are accepted; their complete human reconstructions are not
claimed by this packet.  The separate supplied \texttt{Moduli}
and inherited HUF/PUN foundations likewise remain.  This is
the strongest complete original-scope conditional section-law
join proved here, and supplies no automatic closure of
\texttt{IsUnf} on unsupplied data or of a selected RH/GRH leaf.
\end{proof}
% END GSL26

\subsection*{Exact source correspondence}

{\small\raggedright
All paths in this paragraph are relative to the pinned repository's
\path{src/Development/AbelianVarietyTheory/} unless explicitly
stated otherwise.  Write $A$ for
\path{Proofs/AbelianSchemeLawProofs/} and $C$ for
\path{Proofs/CoherentCohomologyProofs/}.  The complete source files,
exact ranges and SHA-256 bindings are provided in this component's
\path{source_manifest.json} and \path{provenance.json}.

\begin{itemize}
 \item GSL1, GSL11--GSL12 and GSL23:
 $A$\path{SecBaseChangeOfPerf.lean}, lines 22--170,
 241--287 and 305--401.  The original complete conditional
 join is \texttt{secBaseChange\_of\_perf}, lines 333--401.
 \item GSL2--GSL4:
 $A$\path{ExistsMulEqSmulOneOfMulVecInjective.lean},
 lines 9--53, 55--156 and 158--181, respectively.
 \item GSL5:
 $C$\path{ExistsMatricesOfProjectiveComplex.lean},
 especially lines 7--52 and 90--170.
 \item GSL6--GSL7:
 $C$\path{ExistsFiniteProjectiveComplexQuasiIso.lean},
 lines 14--91, 124--244, 255--348 and 521--624.
 \item GSL8:
 $C$\path{ExistsFiniteProjectiveComplexOfNoetherianModel.lean},
 lines 25--60.  The geometric Noetherian-model producer retained
 above is $C$\path{DehomogenisingCompOfSaturating.lean},
 lines 44--114; it is not counted as proved here.
 \item GSL9--GSL10:
 $C$\path{ThreeTermComplexOfProjectiveCechModel.lean},
 lines 24--156; the common-chart and model inputs also appear in
 $C$\path{VeryAmpleCompactQuasiSeparatedAffineCharts.lean},
 lines 24--50, and
 $C$\path{ExistsFiniteProjectiveCechComplexLineBundle.lean},
 lines 19--43.  The last geometric producer remains a premise.
 \item GSL13, GSL14 and GSL22--GSL25:
 $A$\path{FinrankOfIsSecBasis.lean}, lines 62--110,
 126--216, 231--310 and 312--421.  In particular
 \texttt{geometric\_fibre} uses the inverse translation,
 and \texttt{secTwist\_of\_perf} includes every residue lift.
 \item The all-cover isomorphism argument of GSL14 is
 $A$\path{FirstCohomologyVanishesOfIsoOfIsGen.lean},
 lines 129--224.  GSL15 is
 $A$\path{VeryAmpleByOfRelAmple.lean}, lines 29--121.
 \item GSL16--GSL17:
 $A$\path{FirstCohomologyVanishesDescent.lean},
 lines 65--119 and 123--228.  The retained Serre input is the
 first conjunct in
 $A$\path{SerreVanishingVeryAmpleStatement.lean}.
 \item GSL18--GSL19:
 $A$\path{FinrankSectionsLineBundlePowEqPow.lean},
 lines 13--123 and 125--260.  GSL20 uses the base-change
 section equivalence in
 $A$\path{SecMInvariantEqBotOrUnivAtGeometricPoint.lean},
 lines 222--243 and 304--423, and
 $A$\path{SecBaseChangeOfPerf.lean}, lines 305--325.
 \item GSL21:
 $A$\path{TwistedPowTranslate.lean}, lines 63--121.
 The precise assembly boundary in GSL26 is recorded in
 $A$\path{LawsOfH0EqOneClosedFields.lean}, lines 64--91.
\end{itemize}

The literal contracts are in
\path{Definitions/AbelianSchemeLaws/FirstCohomologyVanishesSpecOfIsLocalRing.lean},
lines 16--88, and, relative to the repository's \texttt{src/},
\path{ReciprocityTheory/ReductiveGroups/ArithmeticGeometry/ExtensionDescent/}:
\path{NonzeroLocus.lean}, lines 15--75;
\path{GlobalSectionsRankAlgClosed.lean}, lines 19--41 and 160--166;
and \path{SecBaseChangeAnySt.lean}, lines 16--58.
These give the actual \texttt{PolAV}, \texttt{Perf}, basis and lifting
definitions used throughout.  No source declaration is counted as a
human proof merely because its formal verification is accepted.
\par}
